\exercisesection{1}

\begin{enumerate}
\def\labelenumi{\arabic{enumi}.}
\item
  Let \(d\) be a rational integer not divisible by a square in \(\mathbf{Z}\) . Show that the elements of the field \(\mathbf{K} = \mathbf{Q}(\sqrt{d})\) which are integral over \(\mathbf{Z}\) are of the form \(a + b\sqrt{d}\) where \(a \in \mathbf{Z}\) , \(b \in \mathbf{Z}\) if \(d - 1 \not\equiv 0 \pmod{4}\) and the elements \((a + b\sqrt{d})/2\) where \(a \in \mathbf{Z}\) , \(b \in \mathbf{Z}\) , \(a\) and \(b\) both even or \(a\) and \(b\) both odd, if \(d - 1 \equiv 0 \pmod{4}\) (cf.~Algebra, Chapter VII, § 1, Exercise 8). Deduce that \(\mathbf{A} = \mathbf{Z}[\sqrt{5}]\) is not integrally closed and give an example of an element of \(\mathbf{Q}[\sqrt{5}]\) which is integral over \(\mathbf{A}\) but whose minimal polynomial over the field of fractions of \(\mathbf{A}\) does not have coefficients in \(\mathbf{A}\) .
\item
  Let K be a commutative field and A the sub-K-algebra of the polynomial algebra \(K[X, Y]\) generated by the monomials \(Y^{k}X^{k+1} (k \geqslant 0)\) . Show that XY is such that A{[}XY{]} is contained in a finitely generated A-module but XY is not integral over A.
\item
  Give an example of an infinite sequence \((\mathbf{K}_{n})\) of extensions of a commutative field K of finite degree over K such that the K-algebra \(\prod_{n} K_{n}\) is not algebraic over K.
\item
  In the matrix ring \(\mathbf{M}_{2}(\mathbf{Q})\) give an example of two elements integral over Z but neither the sum nor the product of which is integral over Z (consider matrices of the form \(I + N\) , where N is nilpotent).
\item
  Let A be a commutative ring, B a commutative ring containing A and x an invertible element of B. Show that every element of \(A[x] \cap A[x^{-1}]\) is integral over A. (If \(y = a_{0} + \cdots + a_{p}x^{-p} = b_{0} + \cdots + b_{q}x^{q}\) , where the \(a_{i}\) and \(b_{j}\) are in A, show that the sub-A-module of B generated by 1, \(x, \ldots, x^{p+q+1}\) is a faithful A{[}y{]}-module.)
\item
  Let A be a commutative ring and B a commutative A-algebra; suppose that the minimal prime ideals of B are finite in number; let \(p_{i}\) ( \(1 \leqslant i \leqslant n\) ) be these ideals (note that this hypothesis holds if B is Noetherian). For an element \(x \in B\) to be integral over A, it is necessary and sufficient that each of its canonical images \(x_{i}\) in \(B/p_{i}\) be integral over A. (If this condition holds, show first that x is integral over the subalgebra of B generated by the nilradical \(N = \bigcap_{i} p_{i}\) of B and use the fact that the elements of N are nilpotent.)
\item
  Give an example of a reduced ring A which is integrally closed in its total ring of fractions and which is not a product of integral domains (take the integral closure of a field K in the infinite product \(\prod_{n} K_{n}\) of suitable algebraic extensions of K).
\item
  Let A be a commutative ring and B a finitely generated A-algebra. Show that there exists a finite A-algebra C which is a free A-module and a surjective A-homomorphism \(C \rightarrow B\) .
\item
  Show that the quotient domain \(\mathbf{Z}[\mathbf{X}] / (\mathbf{X}^2 +4)\) is not integrally closed even though \(\mathbf{Z}[\mathbf{X}]\) is.
\item
  Let A be a completely integrally closed domain. Show that for every set I the polynomial domain \(A[X_{i}]_{i \in I}\) and the domain of formal power series \(A[[X_{i}]]_{i \in I}\) (Algebra, Chapter IV, §5, Exercise 1) are completely integrally closed. (Use Proposition 4 of no. 4 and the following lemma: for every finite subset J of I and all \(P \in A[[X_{i}]]_{i \in J}\) , let \(P_{J}\) be the formal power series in \(A[[X_{i}]]_{i \in J}\) derived from P by replacing by 0 in P all the \(X_{i}\) of index \(i \notin J\) ; if P, Q are two formal power series such that \(P_{J}\) is divisible by \(Q_{J}\) in \(A[[X_{i}]]_{i \in J}\) for all J, then P is divisible by Q in \(A[[X_{i}]]_{i \in I}\) . For this note that, if J, \(J'\) are two finite subsets of I such that \(J \subset J'\) and \(P_{J'} = L'Q_{J'}\) and \(P_{J} = LQ_{J}\) , then L is obtained from \(L'\) by replacing by 0 in \(L'\) the \(X_{i}\) of index \(i \in J' - J\) .)
\item
  \begin{enumerate}
  \def\labelenumii{(\alph{enumii})}
  \tightlist
  \item
    Let R be an integral domain, \((\mathbf{A}_{\alpha})\) a right directed family of subrings of R and A the union of the \(A_{\alpha}\) . Show that, if each of the \(A_{\alpha}\) is integrally closed, so is A.
  \end{enumerate}
\end{enumerate}

\begin{enumerate}
\def\labelenumi{(\alph{enumi})}
\setcounter{enumi}{1}
\tightlist
\item
  Let K be a commutative field and \(\mathbf{R} = \mathbf{K}(\mathbf{X}, \mathbf{Y})\) the field of rational functions in two indeterminates over K. For every integer n \textgreater{} 0 let \(A_{n}\) denote the subring \(\mathbf{K}[\mathbf{X}, \mathbf{Y}, \mathbf{Y}\mathbf{X}^{-n}]\) of R. Show that the \(A_{n}\) are completely integrally closed but that their union A is not. (For an element of the form \(\mathbf{P}(\mathbf{X}, \mathbf{Y}) \cdot \mathbf{X}^{-ns}\) , where \(P \in K[X, Y]\) , to belong to \(A_{n}\) show that it is necessary and sufficient that P belong to \(S^{s}\) , where S denotes the ideal of \(K[X, Y]\) generated by \(X^{n}\) and Y.)
\end{enumerate}

¶ * 12. Let A be the ring of integral functions (with complex values) of a complex variable. It is an integral domain whose field of fractions is the field of meromorphic functions of a complex variable.

\begin{enumerate}
\def\labelenumi{(\alph{enumi})}
\item
  Show that \(\mathbf{A}\) is completely integrally closed.
\item
  For all \(x \in \mathbf{A}\) , let \(Z(x)\) be the set of zeros of \(x\) (which is a closed subset of \(\mathbf{C}\) , which is discrete (and therefore countable) if \(x \neq 0\) ). If \(\mathfrak{a}\) is an ideal of \(\mathbf{A}\) distinct from \(\mathbf{A}\) and (0), show that the \(Z(x)\) for \(x \in \mathfrak{a}\) form a filter base \(\mathfrak{F}_{\mathfrak{a}}\) . (If \(x\) and \(y\) are two integral functions with no zero in common, show, using the Mittag-Leffler Theorem, that it is possible to write \(1 / xy = u / x + v / y\) , where \(u\) and \(v\) are in \(\mathbf{A}\) , and deduce that \(Ax + Ay = A\) ). Conversely, for every filter \(\mathfrak{F}\) on \(\mathbf{C}\) to which a closed discrete subset of \(\mathbf{C}\) belongs, the set \(\mathfrak{a}(\mathfrak{F})\) of \(x \in \mathbf{A}\) for which \(Z(x) \in \mathfrak{F}\) is an ideal of \(\mathbf{A}\) . Show that \(\mathfrak{F}_{\mathfrak{a}(\mathfrak{F})} = \mathfrak{F}\) for every filter \(\mathfrak{F}\) to which a closed discrete subset of \(\mathbf{C}\) belongs and \(\mathfrak{a}(\mathfrak{F}) \supset \mathfrak{a}\) for every ideal \(\mathfrak{a}\) of \(\mathbf{A}\) distinct from \(\mathbf{A}\) and (0).
\item
  If \(\mathfrak{p} \neq (0)\) is a prime ideal of \(\mathbf{A}\) , show that \(\mathfrak{F}_{\mathfrak{p}}\) is an ultrafilter. (Restricting it to the case where all the sets of \(\mathfrak{F}_{\mathfrak{p}}\) are infinite, show that, if \(\mathbf{M} \in \mathfrak{F}_{\mathfrak{p}}\) is the union of two infinite sets \(\mathbf{M}'\) , \(\mathbf{M}''\) with no point in common, one of these sets belongs to \(\mathfrak{F}_{\mathfrak{p}}\) .) Conversely, for every ultrafilter \(\mathfrak{U}\) on \(\mathbf{C}\) to which a closed discrete subset of \(\mathbf{C}\) belongs, \(\mathfrak{a}(\mathfrak{U})\) is a maximal ideal of \(\mathbf{A}\) ; obtain the converse.
\item
  Let \(\mathfrak{U}\) be an ultrafilter on \(\mathbf{C}\) , the image under the canonical injection \(\mathbf{Z} \to \mathbf{C}\) of a non-trivial ultrafilter on \(\mathbf{Z}\) . Show that there exist non-maximal prime ideals \(\mathfrak{p}\) of \(\mathbf{A}\) such that \(\mathfrak{F}_{\mathfrak{p}} = \mathfrak{U}\) . (For all \(x \in \mathfrak{a}(\mathfrak{U})\) and all \(n \in \mathbf{Z}\) , let \(\omega_x(n)\) be the order of \(x\) at the point \(n\) ( \(\omega_x(n) = 0\) if \(x(n) \neq 0\) ); consider the set \(\mathfrak{p}\) of \(x \in \mathfrak{a}(\mathfrak{U})\) such that \(\lim_{\mathfrak{u}} \omega_x = +\infty\) .)
\item
  With the notation of (d), let m be the maximal ideal \(\mathfrak{a}(\mathfrak{U})\) . Show that the domain \(A_{m}\) is not completely integrally closed (consider the meromorphic function \(1/(\sin\pi z)\) ).
\end{enumerate}

¶ 13. Let A be a local integral domain and K its field of fractions. Consider the two following properties:

\begin{enumerate}
\def\labelenumi{(\roman{enumi})}
\item
  A is Henselian (Chapter III, § 4, Exercise 3).
\item
  For every finite extension \(\mathbf{L}\) of \(\mathbf{K}\) , the integral closure \(\mathbf{A}'\) of \(\mathbf{A}\) in \(\mathbf{L}\) is a local ring.
\end{enumerate}

Show that (i) implies (ii) and that the converse is true if A is integrally closed. (To see that (i) implies (ii), note that A' is the union of a right directed family of finite A-algebras and use the definition of a Henselian ring. To see that (ii) implies (i) if A is integrally closed, show that, if P ∈ A{[}X{]} is an irreducible monic polynomial and f: A → k is the canonical homomorphism of A onto its residue field k, \(\bar{f}(P)\) is irreducible in K{[}X{]} and consider the field L = K{[}X{]}/(f).)

If condition (ii) is fulfilled, the integral closure A' of A in any extension of K is a local ring. Consider the case of complete Hausdorff local rings; if A is a complete Hausdorff Noetherian local ring, every A-algebra contained in A' and which is a finitely generated A-module is a complete Hausdorff local ring.

\begin{enumerate}
\def\labelenumi{\arabic{enumi}.}
\setcounter{enumi}{13}
\tightlist
\item
  Let A be a completely integrally closed domain and K its field of fractions. Show that for every extension L of K the integral closure A' of A in
\end{enumerate}

L is a completely integrally closed domain. (Reduce it to the case where L is a finite quasi-Galois extension of K of degree m: if \(x \in L\) is such that there exists \(d \in A'\) such that \(dx^{n} \in A'\) for all \(n \geqslant 0\) , show first that it may be assumed that \(d \in A\) , then deduce that, if \(c_{i}\) ( \(1 \leqslant i \leqslant m\) ) are the coefficients of the minimal polynomial of x over K, then \(d^{m}c_{i}^{n} \in A\) for all \(n \geqslant 0\) .)

\begin{enumerate}
\def\labelenumi{\arabic{enumi}.}
\setcounter{enumi}{14}
\tightlist
\item
  Let K be a commutative field of characteristic 0 containing all the roots of unity and such that there exist Galois extensions of K whose Galois group is not solvable (cf.~for example Algebra, Chapter V, Appendix I, no. 1, Proposition 2). Let \(\Omega\) be an algebraic closure of K and let \(\Omega_0\) be the set of elements \(z \in \Omega\) such that the Galois extension of K generated by \(z\) in \(\Omega\) admits a solvable Galois group.
\end{enumerate}

\begin{enumerate}
\def\labelenumi{(\alph{enumi})}
\item
  Show that \(\Omega_0\) is a subgroup of \(\Omega\) distinct from \(\Omega\) (cf.~Algebra, Chapter V, § 10, no. 4, Theorem 1).
\item
  Let A be the subring of \(\Omega[X]\) consisting of the polynomials P such that \(\mathrm{P}(0) \in \Omega_{0}\) . Show that A is not integrally closed but that every element of its field of fractions \(\Omega(X)\) a power of which belongs to A is itself in A.
\end{enumerate}

\begin{enumerate}
\def\labelenumi{\arabic{enumi}.}
\setcounter{enumi}{15}
\item
  Let B be a ring, A a subring of B and f the ideal of A the annihilator of the A-module B/A. Let U be the complement in \(\mathbf{X} = \operatorname{Spec}(\mathbf{A})\) of the closed set \(\mathbf{V}(\mathfrak{f})\) and let \(u = {}^{a}\rho\) , where \(\rho: A \to B\) is the canonical injection. Show that the restriction of u to \({}^{-1}(U)\) is a homeomorphism of \({}^{-1}(U)\) onto U (for every element \(g \in f\) , consider the spectrum of the ring \(A_{g}\) identified with the open set \(X_{g} \subset U\) ).
\item
  Let K be a field, B the polynomial ring \(K[X_{n}]_{n\in N}\) and A the subring of B generated by the monomials \(X_{n}^{2}\) and \(X_{n}^{3}\) for all \(n\in N\) ; show that B is the integral closure of A and that the conductor of B in A is reduced to 0; but, if \(S=A-\{0\}\) , the conductor of \(S^{-1}B\) in \(S^{-1}A\) is not reduced to 0 and \(S^{-1}A\) is integrally closed.
\item
  Let A be an integral domain, K its field of fractions, M a finitely generated A-module, u an endomorphism of M and \(u \otimes 1\) the corresponding endomorphism of the finite dimensional vector K-space \(M_{(K)} = M \otimes_{A} K\) . Show that the coefficients of the characteristic polynomial of \(u \otimes 1\) are integral over A (use condition \((E_{\mathrm{III}})\) of Theorem 1 of no. 1).
\item
  \begin{enumerate}
  \def\labelenumii{(\alph{enumii})}
  \tightlist
  \item
    Let A be an integrally closed domain, K its field of fractions, L a separable algebraic extension of K, x an element of L integral over A, \(K'\) the field \(\mathrm{K}(x) \subset \mathrm{L}\) and F the minimal polynomial of x over K. If x is of degree n over K, show that the integral closure of A in \(K'\) is contained in the A-module generated by the elements \(1/\mathrm{F}'(x)\) , \(x/\mathrm{F}'(x)\) , \(\ldots\) , \(x^{n-1}/\mathrm{F}'(x)\) of \(K'\) . (If \(z \in K'\) is integral over A, write
  \end{enumerate}
\end{enumerate}

\[
z \mathrm{F} ^ {\prime} (x) = g (x), \quad \text { where } \quad g (\mathbf {X}) = b _ {0} + b _ {1} \mathbf {X} + \dots + b _ {n - 1} \mathbf {X} ^ {n - 1} \in \mathbf {K} [ \mathbf {X} ].
\]

Show with the aid of Lagrange's interpolation formula that

\[
g (\mathbf {X}) = \operatorname{Tr} _ {\mathbf {K} ^ {\prime} [ \mathbf {X} ] / \mathbf {K} [ \mathbf {X} ]} \left(z. \frac {\mathrm{F} (\mathbf {X})}{\mathbf {X} - x}\right) \cdot)
\]

\begin{enumerate}
\def\labelenumi{(\alph{enumi})}
\setcounter{enumi}{1}
\tightlist
\item
  Let A be an integrally closed domain, K its field of fractions and p the characteristic exponent of K. Suppose that the subring \(A^{1/p}\) of \(K^{1/p}\) consisting of the p-th roots of the elements of A is a finitely generated A-module. Show that for every finite extension E of K the integral closure of A in E is a finitely generated A-module (reduce to the case where E is a quasi-Galois extension).
\end{enumerate}

¶ * 20. Let \(\mathbf{K}_0\) be a perfect field of characteristic \(p > 0\) , \(\mathbf{K}\) the field of rational functions \(\mathrm{K}_0(\mathrm{X}_n)_{n \in \mathbb{N}}\) and \(\mathbf{B}\) the ring of formal power series \(\mathrm{K}[[\mathrm{Z}]]\) , where \(\mathbf{Z}\) is an indeterminate; if \(m\) is the maximal ideal of \(B\), \(B\) is Hausdorff and complete with the \(m\) -adic topology and is a discrete valuation ring. Let \(L = K((Z))\) be the field of fractions of \(B\) and \(E_0\) the subfield of \(L\) generated by \(L^p\) and the elements \(Z\) and \(X_n\) (for \(n \in N\) ).

\begin{enumerate}
\def\labelenumi{(\alph{enumi})}
\item
  Show that the element \(c = \sum_{n=0}^{\infty} \mathbf{X}_n \mathbf{Z}^n\) of \(L\) does not belong to \(E_0\) .
\item
  Let E be a maximal element of the set of subfields of L containing \(E_{0}\) and not containing c; we write \(A = B \cap E\) . Show that A is a discrete valuation ring (and therefore Noetherian) whose field of fractions is E; if \(m'\) is the maximal ideal of A, then \(m'^{k} = m^{k} \cap A\) and A is dense in B with the m-adic topology; \(L = K(c)\) is an extension of K of degree p and B is the integral closure of A in L but B is not a finitely generated A-module (use Chapter III, §3, Exercise 9 (b)). *
\end{enumerate}

\begin{enumerate}
\def\labelenumi{\arabic{enumi}.}
\setcounter{enumi}{20}
\tightlist
\item
  \begin{enumerate}
  \def\labelenumii{(\alph{enumii})}
  \tightlist
  \item
    Let \(K_{0}\) be a perfect field of characteristic p \textgreater{} 0, L the field of rational functions \(\mathrm{K}_{0}(\mathrm{X}_{n})_{n\in\mathbf{N}}\) , \(A'\) the ring of formal power series \(L[[Y,Z]]\) , where Y and Z are two indeterminates, and A the tensor product ring \(K[[Y,Z]]\otimes_{K}L\) , where we write \(K = L^{p} \subset L\) ; A is a Noetherian local ring, which is not complete and whose completion is identified with \(A'\) (Chapter III, §3, Exercise 17) and \(A'\) is integral over A. Let \(c_{n}\) denote the element
  \end{enumerate}
\end{enumerate}

\[
\mathrm{X} _ {n} \mathrm{Y} + \mathrm{X} _ {n + 1} \mathrm{YZ} + \mathrm{X} _ {n + 2} \mathrm{YZ} ^ {2} + \dots
\]

of A'. Let B be the subring of A' generated by A and the elements \(c_{n}\) ( \(n \geqslant 0\) ). If C = A{[}c₀{]}, C is Noetherian and the ring B (which is not integrally closed) is contained in the integral closure of C and contains C; but show that B is not Noetherian (observe that \(c_{n}\) does not belong to the ideal of B generated by the \(c_{i}\) of index i \textless{} n).

\begin{enumerate}
\def\labelenumi{(\alph{enumi})}
\setcounter{enumi}{1}
\tightlist
\item
  Suppose \(p = 2\) ; \(K_0\) , \(K\) and \(L\) being used with the same meaning as in (a), consider this time the ring of formal power series \(A' = L[[Y, Z, T]]\) in three indeterminates, the subring
\end{enumerate}

\[
\mathrm{A} = \mathrm{K} [ [ \mathrm{Y}, \mathrm{Z}, \mathrm{T} ] ] \otimes_ {\mathrm{K}} \mathrm{L},
\]

and write

\[
b = \mathrm{Y} \sum_ {n \geqslant 0} \mathrm{X} _ {2 n} \mathrm{T} ^ {n} + \mathrm{Z} \sum_ {n \geqslant 0} \mathrm{X} _ {2 n + 1} \mathrm{T} ^ {n}.
\]

Show that the ring A{[}b{]} is Noetherian but that its integral closure B is not a Noetherian ring. (As \(b^{2} \in A\) , every element of the field of fractions of A{[}b{]} is of the form \(P + Qb\) , where P and Q are linear combinations of a finite number of \(X_{k}\) with coefficients in the field of fractions of K{[}{[}Y, Z, T{]}{]}; for such an element to be integral over A{[}b{]}, show that it is necessary and sufficient that it belong to A'. Show that B is generated by the elements

\[
b _ {i} = \mathrm{Y} \sum_ {n \geqslant 0} \mathrm{X} _ {2 (i + n)} \mathrm{T} ^ {n} + \mathrm{Z} \sum_ {n \geqslant 0} \mathrm{X} _ {2 (i + n) + 1} \mathrm{T} ^ {n} \quad \text { for } \quad i \geqslant 0
\]

and conclude the argument as in (a).)

\begin{enumerate}
\def\labelenumi{\arabic{enumi}.}
\setcounter{enumi}{21}
\item
  Let A be an integrally closed domain and K its field of fractions; show that for every extension L of K there exists an integrally closed subdomain B of L whose field of fractions is L and such that \(B \cap K = A\) (consider L as an algebraic extension of a pure transcendental extension \(E = K(X_t)_{t \in I}\) of K and thus reduce it to the case where L is an algebraic extension of K, using Exercise 11 (a)).
\item
  \begin{enumerate}
  \def\labelenumii{(\alph{enumii})}
  \tightlist
  \item
    Let K be a commutative field, n an integer prime to the characteristic of K and P a polynomial in K{[}X{]} with only simple roots (in an algebraic closure of K). If we write \(f(\mathbf{X}, \mathbf{Y}) = \mathbf{Y}^{n} - \mathbf{P}(\mathbf{X})\) , show that, if K contains the n-th roots of unity, the domain \(\mathbf{A} = \mathbf{K}[\mathbf{X}, \mathbf{Y}] / (f)\) is integrally closed and has in its field of fractions the separable extension L of K(X) generated by a root y of f (considered as a polynomial in K(X){[}Y{]}). (Show that, if an element z of L is integral over K{[}X{]}, it belongs to A; for this, write z in the form
  \end{enumerate}
\end{enumerate}

\[
z = \sum_ {i = 0} ^ {n - 1} a _ {i} (\mathbf {X}) y ^ {i},
\]

where \(a_{i}(\mathbf{X}) \in \mathbf{K}(\mathbf{X})\) . Show that the elements \(a_{i}(\mathbf{X})y^{i} (0 \leqslant i \leqslant n - 1)\) are integral over \(\mathbf{K}[\mathbf{X}]\) and deduce that the \((a_{i}(\mathbf{X}))^{n}(\mathbf{P}(\mathbf{X}))^{i}\) are elements of \(\mathbf{K}[\mathbf{X}]\) ; conclude that so are the \(a_{i}(\mathbf{X})\) .

\begin{enumerate}
\def\labelenumi{(\alph{enumi})}
\setcounter{enumi}{1}
\tightlist
\item
  Suppose that K is an imperfect field of characteristic \(p \neq 2\) ; let \(a \in K\) be an element not belonging to \(K^{p}\) , let E be the field \(\mathrm{K}(a^{1/p})\) and let us write \(\mathrm{P}(\mathrm{X}) = \mathrm{X}^{p} - a, f(\mathrm{X}, \mathrm{Y}) = \mathrm{Y}^{2} - \mathrm{P}(\mathrm{X})\) . If A is the integrally closed domain \(\mathrm{K}[\mathrm{X}, \mathrm{Y}]/(f)\) , show that \(B = E \otimes_{K} A\) is an integral domain but not integrally closed (consider the element \(y/(\mathrm{X} - a^{1/p})\) of the field of fractions of B).
\end{enumerate}

\begin{enumerate}
\def\labelenumi{\arabic{enumi}.}
\setcounter{enumi}{23}
\tightlist
\item
  Let \(\Delta\) be a torsion-free commutative group written additively. Let A be a commutative ring, B a commutative A-algebra and \(h\colon\mathbf{A}\to\mathbf{B}\) the ring homomorphism defining the algebra structure on B. Then \(h^{(\Delta)}\colon\mathbf{A}^{(\Delta)}\to\mathbf{B}^{(\Delta)}\) (Algebra,
\end{enumerate}

Chapter II, § 11, Exercise 1) is a (ring) homomorphism of the algebra \(\mathbf{A}^{(\Delta)}\) of the group \(\Delta\) with respect to \(\mathbf{A}\) to the algebra \(\mathbf{B}^{(\Delta)}\) of the group \(\Delta\) with respect to \(\mathbf{B}\) . Let \(\mathbf{P} = \sum_{\alpha \in \Delta} b_{\alpha} \otimes \mathbf{X}^{\alpha}\) be an element of \(\mathbf{B}^{(\Delta)}\) . Show that for \(\mathbf{P}\) to be integral over \(\mathbf{A}^{(\Delta)}\) , it is necessary and sufficient that each of the elements \(b_{\alpha} \in \mathbf{B}\) is integral over \(\mathbf{A}\) . (Reduce it to the case where \(\Delta\) is finitely generated and then to the case \(\Delta = \mathbf{Z}\) .) Show that if \(\mathbf{A}\) is an integral domain and integrally closed so is \(\mathbf{A}^{(\Delta)}\) .

\begin{enumerate}
\def\labelenumi{\arabic{enumi}.}
\setcounter{enumi}{24}
\tightlist
\item
  Let \(\Delta\) be a totally ordered commutative group. Extend Propositions 20 and 21 of no. 8 to graded rings of type \(\Delta\) (for Proposition 20 use Exercise 24; for Proposition 21 generalize Lemma 4 of no. 8 to the case where \(\Delta\) is finitely generated and then prove Proposition 21 by taking the direct limit).
\end{enumerate}

¶ 26. (a) Let A be an integral domain such that for all \(a \neq 0\) in A, the ideal \(Aa\) is the intersection of a family \((q_t)\) of saturations of \(Aa\) with respect to prime ideals \(p_t\) which are minimal among those which contain \(a\) , such that the domains \(A_{p_t}\) are integrally closed (resp. completely integrally closed). Then show that A is integrally closed (resp. completely integrally closed).

\begin{enumerate}
\def\labelenumi{(\alph{enumi})}
\setcounter{enumi}{1}
\item
  Let A be a strongly Laskerian local ring (Chapter IV, § 2, Exercise 28) and m its maximal ideal. Suppose that, for some \(a \in \mathbf{A}\) , m is an immersed prime ideal weakly associated with A/Aa (Chapter IV, § 1, Exercise 17). Show that Aa: m is contained in every ideal which is primary for every prime ideal p ≠ m weakly associated with A/Aa (consider the transporter q: m for such a primary ideal q, cf.~Chapter IV, § 2, Exercise 30). If further A is assumed to be an integral domain, show that the relation p.~(Aa: m) = Aa, for a prime ideal p minimal among those which contain Aa, is impossible (if q is the saturation of A with respect to p, note that the above relation would imply that qA \(_{p}\) = qpA \(_{p}\) in A \(_{p}\) and conclude with the aid of Exercise 29 (d) of Chapter IV, § 2).
\item
  Show that, if A is a strongly Laskerian integral domain and there exists an element \(a \neq 0\) of A such that there are prime ideals which are weakly associated with A/Aa and immersed, A is not completely integrally closed. (Reduce it to the case considered in (b); in the notation of (b), there then exists \(b \in \mathbf{A}a\) : m such that \(b \notin \mathbf{A}a\) and \(bp \subset am\) ; deduce by induction on \(n\) that \(b^n p \subset a^n m\) for all \(n\) .)
\item
  Deduce in particular from (a) and (c) that, for a Noetherian integral domain to be integrally closed, it is necessary and sufficient that for all \(a \neq 0\) in A the prime ideals \(\mathfrak{p}\) associated with A/Aa be not immersed and be such that \(\mathrm{A}_{\mathfrak{p}}\) is integrally closed (cf.~Chapter VII, § 1, no. 4, Proposition 8).
\end{enumerate}

\begin{enumerate}
\def\labelenumi{\arabic{enumi}.}
\setcounter{enumi}{26}
\tightlist
\item
  Let A be an integrally closed but not completely integrally closed domain and K its field of fractions; let \(a \neq 0\) be a non-invertible element of A for which there exists \(d \neq 0\) in A such that \(da^{-n} \in A\) for every integer \(n \geqslant 0\) .
\end{enumerate}

Show that in the field of formal power series \(\mathbf{K}((\mathbf{X}))\) there exists a formal power series \(\sum_{k=1}^{\infty} u_k \mathbf{X}^k = f(\mathbf{X})\) satisfying the equation

\[
(f (\mathbf {X})) ^ {2} - a f (\mathbf {X}) + \mathbf {X} = 0,
\]

hence integral over A{[}{[}X{]}{]}, belonging to the field of fractions of A{[}{[}X{]}{]} but not belonging to A{[}{[}X{]}{]}, so that A{[}{[}X{]}{]} is not integrally closed.

\begin{enumerate}
\def\labelenumi{\arabic{enumi}.}
\setcounter{enumi}{27}
\item
  Let \(k\) be a field and \(\mathbf{A}_0\) the polynomial ring \(k[\mathrm{X}_n, \mathrm{Y}_n]_{n \in \mathbf{Z}}\) in two infinite families of indeterminates. Let \(G\) be an infinite monogenous group (therefore isomorphic to \(\mathbf{Z}\) ) and let \(\sigma\) be a generator of \(\mathcal{G}\) . \(G\) is defined as operating on \(\mathbf{A}_0\) by the conditions \(\sigma(\mathrm{Y}_n) = \mathrm{Y}_n\) and \(\sigma(\mathrm{X}_n) = \mathrm{X}_{n+1}\) for all \(n \in \mathbf{Z}\) . Let \(\Im\) be the ideal of \(\mathbf{A}_0\) generated by the elements \(\mathrm{Y}_n(\mathrm{X}_n - \mathrm{X}_{n+1})\) for all \(n \in \mathbf{Z}\) and let \(A\) be the quotient ring \(A_0 / \Im\) ; as \(\sigma(\Im) \subset \Im\) , the group \(\mathcal{G}\) also operates on \(A\) by passing to the quotient. Let \(S\) be the multiplicative subset of \(A\) generated by the canonical images of the \(Y_n\) in \(A\) , so that \(S\) consists of the elements invariant under \(G\) . Show that \(S^{-1}(A^{\mathcal{G}}) \neq (S^{-1}A)^{\mathcal{G}}\) .
\item
  Let k be a field of characteristic \(\neq2\) , A the polynomial ring k{[}T{]} in one indeterminate and a, b two elements of k such that \(a\neq0\) , \(b\neq0\) and \(a\neq b\) . Let \(\mathbf{K}=k(\mathbf{T})\) be the field of fractions of A and L, M the extensions of K adjoining to K respectively a root of \(\mathbf{X}^{2}-\mathbf{T}(\mathbf{T}-a)\) and a root of \(\mathbf{X}^{2}-\mathbf{T}(\mathbf{T}-b)\) . Let B, C be the integral closures of A in L and M respectively, which are finitely generated free A-modules. Show that \(L\otimes_{K}M\) is a field, a separable finite extension of K, but that \(B\otimes_{A}C\) , which is canonically identified with a subring of \(L\otimes_{K}M\) , is not the integral closure of A in \(L\otimes_{K}M\) .
\end{enumerate}
% semantic-fragments: legacy-input-seam
\relax{}
